\emph{By Hal Draper}

For social change toward a better world, socialists believe the most important and indeed decisive social force is the struggle of the working class. Why the working class?

Why do socialists believe there is a special connection between their own great goal of a new society and the interests of labor, this one segment of society? Is it because we "idealize" workers as being better, or more clever, or more honest, or more courageous, or more humanitarian, than non-workers?

\emph{--Isn't it rather true that the workers have time and again followed reactionary courses and leaders and have by no means shown any invariable affinity for progressive causes? Don't they follow Stalinist totalitarianism in countries like France and Italy; and where they do not, are their own trade-union bureaucrats much of an asset? Haven't they been misled and deceived like any other section of society? Aren't they filled with race prejudice in the U.S., sometimes even more so than the upper classes? If it is true that workers are "naturally" pro-socialist, why is it they have made such a mess of things, voting for reactionaries and fakers and supporting the status quo?} ... And so on.

Most of this type of questioning is based on simple misunderstanding of the socialist viewpoint about the working class. Especially in this country, where the socialist movement has always been relatively weak, the most popular anti-socialist notions are mostt often founded on misinformation about what socialists believe, because their voices have not been loudly heard.

Socialists do not "idealize" workers in any sense whatever.

Taking them man for man, as individuals, there is no reason to argue whether workers are "better" human beings than others because they are workers. This whole approach, whether pro or con, has nothing to do with the socialist conception.

To underline this in a different way: If we try to view social issues as mere conflicts between Good People and Bad People, then surely we must say that men who insist on starving others are Bad. The present minimum wage is surely a pittance; yet opposition even to t his pittance was strong among employers, especially small employers, while virtually absent among workers. Is this a tendency of employers because they are Bad Men? On the contrary, these employers are just as likely to be kind fathers, generous friends, indulgent husbands, charity-givers--not the type to deliberately run over children in the street. They act one way as individual atoms in the social fabric; they act in quite another way as part of their class collectivity.

They explain this, when they do, by saying "Business is business." This is their way of distinguishing their individual and human thoughts and role from their role as a member of the business community--that is, of their class. In the latter case, the conditions of existence and interests of "business" make out of them a social force that has little resemblance to their individual psychologies.

Like every other class or group, the working class is more than the sum of its individual atoms.

Man for man, workers are not "naturally" more pro-socialist than anyone else. It is a question of what \emph{direction} they are pushed in by the conditions of their existence \emph{as a class} and by their interests \emph{as workers}, just as with any other group.

This indeed is one reason why so often socialist ideas tend to be \emph{initiated} in a systematic way not by ideologists from the working class but by men from t he "educated classes," the bourgeoisie and intellectuals, men like Marx and Engels, for example, who were not proletarians themselves--altho it should be noted that the impulsions to the systematization of such ideas were coming from the working masses' struggles and conditions, not from other sections of society. Individuals were led to align themselves with the working class.

If they were drawn in this direction, it was because here was the dynamic social force which they recognized as the decisive one for putting flesh and blood on ideas.

When a working class is politically and socially undeveloped it is well-nigh inevitable that its members will be filled with all sorts of backward and even reactionary notions. For example, it has often been found in the U.S. that racial intolerance decreases with amount of education: college graduates are less prejudiced, \&c. Now, in general, working-class children get less schooling than upper-class offspring. So according to this pattern, workers should be far more filled with racism than the middle class. It is instructive to see where this neat pattern does and \emph{does not} hold.

It holds best where labor is most poorly organized as a class, and most recently organized, and where it is oprganized in the least class-conscious fashion. The South is not only a cauldron of racism but also a sinkhole of union-busting and open-shopism. Toward the other end of the scale, racism is combated--as no-where in middle-class groups--in the more militant mass-production unions that sprang from the CIO upheaval, like the Auto Workers, not to speak of the socialist movement.

Here anti-racism is not a function of school education; it is a function of class education. In many a mass-production integrated local, the organization is often more anti-racist than the sum of its members. That is, the dynamics of \emph{class} needs push it more strongly against racism, which is divisive of the class, than do the individual opinions of its members.

What we have been emphasizing is that the socialist sees no special magic in the "worker" as an atomized individual. The special "advantage" of the working class springs from inherent drives of its \emph{class} position in society, its ineradicable interests as a group, its conditions of life; and this "advantage" comes into play insofar as this class organizes itself (as it is inevitably driven to do so) and transforms the thinking of its individual components in the course of class experiences.

Now it is this sort of thing that the socialist calls the development of class-consciousness. The U.S. is the one modern country in the world where the working class is still at a rather elementary stage of class-consciousness. Therefore it is particularly in this country, and most particularly among academicians who have no roots in the real social struggle of our times, that the special role of the working class is most persistently questioned.

It is much harder to do so in Great Britain, for example, where this "special role of the working class" is the daily headache of the Tories, and also of the Labor Party leaders themselves. Or in France and Italy, where the strength of the Communist Parties is closely connected precisely with their ability to use and abuse "the special role of the working class." Or in almost any other European country, where the working class is well organized as a class. Or even in leading countries of Asia and Latin America, where working-class-based forces play prominent roles out of all proportion to the size of the class.

In this respect it is the U.S. which is "out of step," which is the exception to the rule; and while American bourgeois ideologists may be grateful for this exceptional position, they have no license to deny the rule.

The "rule" is that all over the world organized working-class struggle is inextricably bound up with every effort toward freedom and human emancipation. Where the working class has been defeated, democracy and progress and humanity have been defeated too. Where the forces of freedom have fought, in Hungary 1956 as in capitalist Europe, it is the working-class forces that have been in the van.

There is no other sector of society of which this or anything like it can be said--not the middle class, not the intellectuals, not the "educated classes," not the students, not the "managers," not anyone elese except the organized working class, for good or ill.

What is the "advantage" which the working class possesses willy-nilly, by virtue of the terms of its own existence under capitalism? Here in outline form are the special characteristics inherent in a social class whose individual human components are (remember) no better or worse than the rest.

(1) The conditions of life of the working class lead it to \emph{organize} in the first place--and most solidly as a homogenous movement.

There is, of course, one other class which rivals the working class in this respect: the capitalists themselves, whose own class-consciousness and sense of class-solidarity are ever-present models for the workers.

Never has a predominantly agrarian population (farmers or peasants) been able to duplicate the organizational achievements of the working class. The difference is no reflection on the individual farmer. By terms of their life, they live in atomized groups which stress self-sufficiency, separateness, reliance on individual effort; they are not thrown together in crowds and subjected to simultaneous stresses in the heat of social struggles as are workers.

Workers are taught organization not by superior intelligence or outside agitators, but by the capitalists themselves. They are organized on the assembly lines, in the factory gangs, in shifts, in work teams, in the division of labor of capitalism itself. Capitalism cannot live without "organizing" its workers, teaching them the virtues of working together, therefore of solidarity.

It teaches discipline. It enforces centralization of effort. It hammers home every day the advantages of pooled work, and the subordination of individual self-interest to the needs of a group.

It does not teach this lesson equally to all workers: it is plainer for assembly-line workers in the mass-production industries than (say) for an office secretary who takes dictation from a personal boss, who works \emph{with} a boss rather than \emph{with} fellow workers. This fact links up also with the social views which arise among these different strata of workers--simply on the basis of this first point: class \emph{organization}.

(2) The interests of workers as a solidarized group, organized by capitalism, lead them to \emph{struggle}.

It must be emphasized that this often takes place quite apart from the conscious desires and wishes of the labor leaders themselves. Labor leaders, risen from the ranks of lowly workers and aspiring to be accepted as respectable and responsible members of bourgeois society, often \emph{want} to substitute pleasant and friendly conferences with management for any kind of conflict. Having freed themselves from the condition of existence to which the mass of workers are condemned, they tend to become "bourgeoisified"; they want to integrate into the ruling class, or at least find as respectable a niche there as a corporation lawyer.

And indeed they could do so (so many do!) if not for the fact that it is the working class that they are standing on in order to reach so high. For the working class needs representatives in order to oppose the bosses' interests; but the bosses accept the friendship of these labor leaders only insofar as they "behave." From \emph{below} these bourgeoisified bureaucrats, there always arises the pressure of mass demands, the unslakable needs of the workers which cannot be wished away with fine talk about class collaboration, the aspirations steaming up from the depths of the class, demanding "delivery of the goods."

Some bureaucrats can continue their precarious balancing-act for substantial periods, in "normal" times of class quiet particularly, as everybody knows; but even the most conservative and most bourgeoisified union leader must to some extent satisfy the class needs of his constituent base, or else--. This is in the worst case, of course, and there are not a few such "worst" cases in the society-corrupted labor bureaucracy of this country. But whether timidly or militantly, consistently or crudely, even the conservative union leader who does not "believe" in class struggle must be its instrument, to the extent that he functions as a labor leader at all.

(3) The direction of the workers' organized struggle inevitably tends to be \emph{counter to capitalism--}or, more finely, this struggle always tends to \emph{go outside the framework of capitalist institutions and ideas.}

Steadily the labor movement's insistence on \emph{social-responsibility} for all aspects of life comes in conflict with the capitalist insistence on the rights of private property. For the essence of capitalist private-property relations is that this whole area of man's life--the economic sphere--is to be withdrawn from the rule of social responsibility, and is to be ruled by the unilateral power of capital as its birthright.

Capitalism has been forced into many compromises in this respect, as is well known--mainly this one, that (a) the state is accorded power to intervene as representative of "society," provided (b) that the associated capitalist class retain full control of this intervening state. (This is the process of "statification" under capitalism in a nutshell.) But whatever the compromises, the working-class movement can never be satisfied--not even the undeveloped union-conscious labor movement of this country.

More militant unions (e.g. UAW) have raised demands like trade-union intervention in the setting of prices or in peering over the capitalists' books to check their profit. In periods of intense class-struggle, sit-downers have taken over the factories without a qualm over the rights of private property. The tendency of the unions in politics is to support social controls all the way down the line--over offshore oil, natural gas, prices, health insurance, etc.--in the name of social responsibility vs. private property. Insofar as this support of "statification" takes place without concomitant insistence on control by a socialized democracy, this is indeed a contribution to the bureaucratization of capitalism rather than its democratization. But given a socialist framework, it is this insistence on social responsibility vs. private property which is the germ of the labor movement's inherent and ineradicable "creeping socialism."

The intuition of the reactionaries is not altogether baseless in this respect, though often exaggerated and viciously directed. Even Samuel Gompers used to argue that his simple slogan of \emph{"More!"} for the labor movement was a more "revolutionary" slogan than the socialists'. At any rate, it is true that. insofar as labor consistently presses for "more" out of the economic pie even when this is incompatible with capitalist needs--insofar as labor presses for "more" social responsibility and less rule by private profit--insofar as labor presses in this direction without drawing back when the capitalists yell too violently--to this extent labor drives the logic of its own existence outside the bounds of the capitalist framework, and tends to explode it.

Of course, we socialists would maintain, and experience shows, that this does not happen except when the working-class movement grows up to adopting socialist leadership and program; but all we are stressing in the present connection is that the class conditions and needs and interests of the workers drive their organized movement, in the course of its struggle, right up against the bounds of the capitalist system.

This is not true of any other group in society--only of individuals from other classes, who may decide to throw in their lot with the working-class struggle. It is enlightening, for example, to study the type of political program commonly adopted by non-working-class parties which set out to express protest: radical peasant parties, or urban middle-class reform parties, or farmers' parties in the U.S.

Peasant parties most typically stop well short of proposing the abolition of capitalism, condining themselves to proposals for improving their class's lot in ways compatible with the rule of private property; for the peasant is a very tenacious private-property holder himself and does not easily see beyond this class limitation. In a different kind of case, as in the Nazy appeal to middle-class elements, a kind of pseudo-anti-capitalism may be patched up by directing slogans against bank capital as distinct from "good" productive capital; or, as in the case of Henry Wallace's program, supporting "progressive" capitalists against "reactionary" capitalists.

But what is noteworthy is this: only in the case of working-class parties, all over the world, does the program and goal of the movement turn fast or slow toward a basic assault on the fundamentals of the capitalist system itself.

Obviously most Americans will not consider that this is a good thing! But the fact itself is what we point to, as illuminating the "special role of the working class," for the benefit of Americans who cannot see that the working class \emph{as a class} does play any special role whatever.

(4) The conditions and interests of the working class not only push it toward organized struggle against capitalism, but impel it toward \emph{a courage and boldness and militancy} which are well-nigh unique to it, at critical moments of struggle when these qualities are called for.

Now at first blush this may seem to be in contradiction with our earlier statements that workers are not necessarily personally "better" in any sense. Are we now saying that workers are braver and bolder, etc.?

Only with the same qualifications previously explained. We are talking about their potentialities as an organized class--plus, perhaps, for many individuals whatever carryover takes place from organized behavior to personal behavior as a result of education in struggle and conditioning in life situations. But it is the class behavior we are interested in.

Stereotypes may be bad, but class "stereotypes" contain more than a kernel of truth. Thus, there is the "Timid Professor." We have known many professors who were not at all personally timid; yet the sweeping stereotype contains a truth about the impact of academic life and its pressures upon the social psychology of professors.

In his \emph{White Collar}, a study of the middle class in America, C. Wright Mills (a non-timid professor) drew a generalized picture of the new middle class which is relevant here. They are the "rearguarders," says Mills, waiting for someone else to move. As a group they have no cohesion, but are on sale to the highest bidder or the most likely winner. "They have no steady discontent or responsible struggle with the conditions of their lives. For discontent of this sort requires imagination, even a little vision; and responsible struggle requires leadership." As individuals with private positions (Mills continues) "they hesitate, confused and vacillating in their opinions, unfocused and discontinuous in their actions... they have no targets on which to focus their worry and distrust. They may be politically irritable, but the have no political passion. They are a chorus, too afraid to grumble, too hysterical in their applause." In the short run, he concludes, they follow the panicky way of prestige; in the long run they follow the ways of power.

This scathing portrait is a picture of a social class, not an insult directed against middle-class individuals, just as we have been discussing the social potentialities of a class and not "idealizing" workers.

But surely, realizing the truth of this portrait, one can see why middle-class groups simply cannot work up the dynamic drive which is necessary before one can be "courageous and bold and militant."

Take a simple model: A factory worker on a picket line can and often does abuse entering scabs and may even have to be restrained from physical attack; he is not constrained by notions of bourgeois respectability, even though he may be quite "respectable" and "bourgeois" on normal occasions. He is, in fact, more alienated from class society, no matter how he thinks, or how he thinks he thinks. But now go along the scale of workers up (or down) toward more and more "respectable" white-collar workers and employees, to office employees, to bank tellers, to fashion-house fitters, to...college professors. And try to imagine them yelling at scabs on a picket line.

We use this example only as a handy and visualizable token of what is involved: the dynamism of the class in its organized struggle for "something better." History provides a better record--the record of the working class in far more crucial situations than mere strikes: records of the heights of valor and self-sacrifice that have been reached by unknown workers, not named heroes, in revolutionary struggles. But these things are not visualizable for the average American, who after all is himself the product of a society dominated by middle-class mediocrity.

(5) Finally, we are talking about the organized and militant anti-capitalist struggle of \emph{the only class which has the social power and weight} to abolish the old order and build a new society.

Whatever a historian may say about the role of force in revolutions, it is a Marxist principle that social revolutions are not made by bullets. This is a caricature of socialist revolution spread by certain types of policement and certain types of professors. The Marxist socialist believes that when the working class, and its associated allies from other sections of the people, are in their massed majority ready for the abolition of capitalism, it is their \emph{social} power which will determine the result in the last analysis.

The social power of the class depends not only on its numbers. It depends also on its homogeneity and organizability, as we have discussed--its striking power. It also depends on t he indispensability of the services which it performs in keeping the society's work going.

No other class has its hands so closely on the basic work without which the system grinds to a halt. Not a wheel can turn without them. No other class can precipitate a social crisis by the deliberate decision of its organized cadres as in a large-scale strike. When the working class goes into battle, all of the society is embroiled, for all depends on it. Everytime the working class stirs, the rest of society quivers. Yet there is debate over its "special role."

After all of the above, there is still a deeper "why" to be asked, a question that goes behind all of the points we have made up to now. Within the confines of this article we can only point to it.

In the last analysis, the "rearguard" character of the middle classes, which Mills pointed to, reflects their political and social blind-alley. They cannot give society a lead because there is no \emph{social program} which effectively corresponds to the special interesats of the middle classes. From the conditions of \emph{their} existence arises no pointer to a way out for all of society.

In contrast, the working class, as the bottom layer of all classes, cannot even stir without \emph{pointing} to a program, even when it itself rejects it: the abolition of capitalism, its class antagonist, and the assumption of social responsibility by the people democratically organized, regardless of private profit.

At bottom, it is because the interests of the working class, implicit in its struggles, point a program for a basic transformation and reconstruction of society, that this class is pushed to take a vanguard role in every struggle for freedom and emancipation.

We need hardly spend much space affirming how cognizant we are of how often the working class and its interests have been deceived and betrayed by its enemies and false friends. The history of capitalism, from one point of view, is nothing but a history of continued duping of the working class. In fact, deception of the working class is one of the most important conditions for the maintenance of capitalism or any other exploitative system.

It is not really necessary for us to learn all about this from critics who like to argue that socialists' "faith" in this working class is misplaced. It is hardley necessary for us to be reminded, also, that today in good part the Communist Parties live by their ability to dupe and deceive the working class in countries like France and Italy. The battle for socialist democracy against both capitalism and Stalinism can even be summed up as the battle to free the working class from its deception by each of these class enemies.

But this is a battle which, by definition, is won as soon as the workers are "undeceived." It is meanwhile a downright irrelevancy in \emph{this} connection for critics to tell us, as they do regularly, that because the working class has so long been deceived and betrayed, we must conclude that it is hopeless.

We point out only this: It is the \emph{working class} that it is crucial for reaction to deceive, not the middle classes or any of the "rearguarders."

The socialist revolution, once observed by Rosa Luxemburg, is a war in which there are necessarily a continuous series of "defeats" followed by only one victory. Nothing can be guaranteed, of course, except the honor and dignity of fighting for a new and better world, rather than the vileness of adapting one's mind and heart to a vile one. We guarantee to no one that the working class is predestined to "behave according to our blueprints" even if \emph{we} sit by in interested passivity to see whether \emph{it} carries out its "mission." We offer only a road of struggle and a choice of allies in the only war worth fighting, the battle for a socialist democracy against the rival world blocs of war and exploitation.
